\section{Why the preserved source is calibrated to PA}
\label{sec:calibration}

This section makes the logical dependency explicit. The detailed supplied
calibration is retained, with provenance, in \file{report/reference/}.
We use its syntactic proof transformations, rather than inferring a PA theorem
from external standard-model nonhalting.

\subsection{The numerical formulas and separate diagonal head}

For Cantor pairing
\[
 p(a,b)=\frac{(a+b)(a+b+1)}2+b,
\]
write
\[
 I(A,B)=4p(A,B),\quad M(a,b)=4p(a,b)+1,\quad
 U(a,B)=4p(B,a)+2,\quad E(a,b)=4p(a,b)+3.
\]
Code zero denotes truth, code one is the contradiction $M(0,0)$, and
$N(B)=I(B,1)$ is negation. The quantifier arguments are transposed as shown.
Conjunction and existential quantification have their usual classical
implicational definitions.

The retained diagonal clause and the guarded old gate are
\begin{align*}
 \Gamma(B)&=B\to U(B,B),\\
 V(a,b,c)&=M(a,b)\to U(c,M(a,b)),\\
 P(a,b,c)&=c\to M(c,b),\\
 S(a,b,c)&=U\bigl(b,(M(a,b)\to P(a,b,c))\wedge(M(c,b)\to c)\bigr),\\
 Q(a,b,c)&=U(c,S(a,b,c)\to N(P(a,b,c))),\\
 G(a,b,c)&=N(V(a,b,c)\to Q(a,b,c)).
\end{align*}
The separate-head AX command pushes $\Gamma(c)$ and leaves $G(a,b,c)$
as its top when $c\ne a,b$; the guarded top is truth otherwise.
The diagonal head is safe without this guard: every variable label in
formula code $B$ is strictly below $B$, so $x_B$ is absent from its own
formula. This follows by induction on the decoding, using strict component
bounds for nonzero constructors and the separate truth-code case.

The other six universally instantiated heads are the PA196 heads:
\begin{align*}
 &(\exists x_a\,U(a,C))\to C,\\
 &U(a,U(b,C))\to U(b,U(a,C)),\\
 &M(c,a)\to(E(c,b)\to M(b,a)),\\
 &E(c,a)\to(M(b,c)\to M(b,a)),\\
 &\exists x_a E(c,a),\\
 &U(a,b)\to\bigl(U(a,b\to c)\to U(a,c)\bigr).
\end{align*}
There is also the unchanged dominated implicational head
$((H\to d)\to D)\to((D\to H)\to(A\to H))$ with the inherited numerical
dominance condition $d\le D$, plus modus ponens and generalization of theorems.
The finite command emitter is checked over a free algebra, including the
order of the pushed formulas and the guard.

The bundled diagonal gate is
\[
 F_\Gamma(a,b,c)=N\bigl(\Gamma(c)\to(V(a,b,c)\to Q(a,b,c))\bigr).
\]
Classical implicational templates convert $F_\Gamma$ into both $\Gamma(c)$
and $G(a,b,c)$, and conversely. Every $\Gamma(B)$ is obtainable from a
bundled gate with $c=B$ and $a=b=B+1$. Thus the bundled and separate-head
calculi have primitive-recursive proof translations preserving conclusions.
The source-level separation changes the stack protocol, so the command
compiler must account for ten pushed entries.

\subsection{Completeness: PA refutations produce source halts}

The supplied PA196 refutation compiler relies on two-way membership
rewriting, the neighborhood-equality implication, and an equality witness.
The diagonal calculus recovers these interfaces without assuming the
removed equality-vacuity schema.

For a reverse rewrite, put
\[
 C=E(a,b)\to(\chi(b)\to\chi(a)),\qquad
 \varphi=N(C),\qquad c=\ulcorner\varphi\urcorner,
\]
where $\chi(t)$ is a membership atom with $t$ in either slot, or $M(t,t)$.
First handle distinct labels, so the code bounds make the displayed
separator label fresh. Pointed separation at $(a,a,c)$ supplies $N(Q)$.
The second half of its separation matrix and the propositional consequence
$\varphi\to N(\chi(a))$ give
\[
 S\to U(a,M(c,a)\to N(\chi(a))).
\]
The retained equality witness, direct membership rewriting and
membership-only vacuity yield
\[
 S\to(M(c,b)\to N(\chi(b))).
\]
Under $\varphi$, the point condition $P$ supplies $M(c,a)$; direct rewriting
supplies $M(c,b)$, whereas $\varphi$ also supplies $\chi(b)$. Discharging
assumptions gives $\varphi\to(S\to N(P))$.
Generalization at $c$ and the available diagonal clause
$\Gamma(\varphi)$ yield $\varphi\to Q$, contradicting $N(Q)$.
Double-negation elimination gives $C$. Coincident-label cases reduce to
propositional reasoning and the already obtained diagonal membership case.

For neighborhood equality, first choose the fresh binder
$z_0=\ulcorner E(a,b)\urcorner$. The instance $\Gamma(E(a,b))$, universal
distribution and the two-way membership rewrites give the neighborhood
implication at $z_0$. The equality witness and membership-only vacuity
transport the binder to each $z\ne a,b$. This also yields the neighborhood
witness. These are the required replacement interfaces in the original
compiler. Code construction, comparisons of labels and finite proof
concatenation are primitive recursive, and PA checks their correctness.

For the ten-entry separate-head command, the adjusted stack words include
\[
\begin{aligned}
 \mathrm{pop}&=\mathsf{AX}(0,0,0)\,\mathsf S^{11},\\
 \mathrm{push}&=\mathsf{AX}(0,0,0)\,\mathsf S^9\,\mathsf{MP}(0),\\
 \mathrm{discard}&=\mathsf{MP}(1)\,\mathsf{AX}(1,0,0)\,
                       \mathsf S^{10}\,\mathsf{MP}(0).
\end{aligned}
\]
The generalization selector uses $\mathsf S^{10}$ after AX. The prior
nine-entry selectors cannot be reused unchanged. The inherited exhaustive
command-word scheduler, applied to this correctly adjusted compiler,
therefore gives a PA-verifiable translation from each PA refutation into
a finite halting source computation.

\subsection{Soundness inside PA: source halts produce refutations}

Interpret $M(a,b)$ by finite binary membership $\operatorname{Bit}(x_a,x_b)$,
$E(a,b)$ by equality, and $U(a,B)$ by universal quantification of the
translated formula. A convenient bit formula is
\[
 \operatorname{Bit}(s,z)\equiv
 \exists u,v\,[s=(2u+1)2^z+v\ \wedge\ v<2^z].
\]
PA proves $\operatorname{Bit}(s,z)\to z<s$, so it proves the negation of
the translation of code one. The exceptional constructor equality
$I(0,0)=0$ is handled by a fixed proof of the equivalence of truth with
truth implying truth; it is not treated as syntactic disjointness of
constructors.

For every arithmetic formula $\theta(z,\vec u)$, PA has a proof of finite
binary comprehension
\[
 \forall n\,\exists s<2^n\,\forall z<n\,
       (\operatorname{Bit}(s,z)\leftrightarrow\theta(z,\vec u)).
\]
The proof is induction on $n$. Retain the previous binary code if
$\neg\theta(n,\vec u)$, and add $2^n$ if $\theta(n,\vec u)$. Substituting
the syntax of $\theta$ into this fixed induction-proof pattern is a
primitive-recursive proof-producing operation. PA verifies that it produces
a proof. This statement concerns syntax and proof checking, not a uniform
truth predicate for arithmetic.

For pointed separation with distinct labels $a,b$, the relevant finite
comprehension matrix is
\[
 (\operatorname{Bit}(x_a,z)\lor z=x_b)\wedge T(c)[z/x_b],
\]
with a suitable bound such as $x_a+x_b+1$. When $a=b$, the binder changes
both occurrences of that label; the correct matrix is instead
$(z=x_b)\wedge T(c)[z/x_b]$. The first antecedent inside the binder is
then $\operatorname{Bit}(z,z)$, which PA refutes. This distinction prevents
a false free-variable treatment of the collision case. Code freshness
provides the required separator label. The diagonal head and the other
logical heads have uniformly generated first-order proof templates.

Annotate every visible theorem code, the cached head, and every live
canonical stack entry with a PA proof of its translation. Both visible
proofs matter because the original Shift halt test is $t+h=1$.
Canonical expansion of the Cantor stack treats $p(0,0)=0$ explicitly:
popping numerical zero gives the fixed truth proof; pushing a zero onto
zero adds no live canonical cell. AX, MP and Shift preserve this annotation
by primitive-recursive proof transformations. A source halt produces a
PA proof of the translated contradiction and hence a PA refutation.
The scheduler carries its store, so this invariant is stated for every
annotated store, without assuming a reset between command words.

This yields the two syntactic implications
\[
 \PA\vdash H_R\ \leftrightarrow\
       \exists p\,\Prf_{\PA}(p,\ulcorner0=1\urcorner).
\]
No soundness or reflection principle for PA is used inside PA.

\begin{theorem}[Consistency calibration of the principal table]
\[
 \boxed{\displaystyle
 \PA\vdash\bigl(\neg H_{\MX}\leftrightarrow\Con(\PA)\bigr).}
\]
\end{theorem}
\begin{proof}
Compose the last displayed source calibration with
Theorem~\ref{thm:transport}, then negate the existential PA-refutation
predicate. The same finite-trace reasoning applies to the register-allocation companion,
with its explicitly documented source-frame conjugacy.
\end{proof}

